\chapter{Poisson's Equation and the Newtonian Potential}

For integrable $f$ on a bounded domain $\Omega$, its Newtonian potential is
\[
w(x)=\int_\Omega\Gamma(x-y)f(y)\,dy,
\]
where $\Gamma$ is the normalized fundamental solution (2.12). Green's
representation decomposes a sufficiently regular function into a harmonic
part and the Newtonian potential of its Laplacian.

\section*{4.1. H\"older Spaces}

\begin{definition}\label{gt:def-holder-pointwise}\dcref{ch4:4.3}\group{ch04-poisson-newtonian-potential}
\source{gilbarg-trudinger:page-0064}
For $0<\alpha<1$, a function $f$ on bounded $D\ni x_0$ is
$\alpha$-H\"older continuous at $x_0$ when
\begin{equation}
\tag{4.3}
[f]_{\alpha;x_0}=\sup_{x\in D,\,x\neq x_0}
\frac{|f(x)-f(x_0)|}{|x-x_0|^\alpha}<\infty.
\end{equation}
For $\alpha=1$ this is pointwise Lipschitz continuity.
\end{definition}

\begin{definition}\label{gt:def-holder-uniform-local}\dcref{ch4:4.4}\group{ch04-poisson-newtonian-potential}
\source{gilbarg-trudinger:page-0064}
A function on $D$ is uniformly $\alpha$-H\"older continuous if
\begin{equation}
\tag{4.4}
[f]_{\alpha;D}=\sup_{x\neq y\in D}
\frac{|f(x)-f(y)|}{|x-y|^\alpha}<\infty.
\end{equation}
It is locally H\"older continuous if this holds on every compact subset.
\end{definition}

\begin{definition}\label{gt:def-holder-spaces}\dcref{ch4:4.5,ch4:4.6,ch4:4.6',ch4:4.7}\group{ch04-poisson-newtonian-potential}
\source{gilbarg-trudinger:page-0065}
The space $C^{k,\alpha}(\overline\Omega)$ consists of
$u\in C^k(\overline\Omega)$ whose derivatives of order $k$ are uniformly
$\alpha$-H\"older. Set
\begin{equation}
\tag{4.5}
\begin{aligned}
[u]_{k,0;\Omega}&=\max_{|\beta|=k}\sup_\Omega|D^\beta u|,
& [u]_{k,\alpha;\Omega}&=\max_{|\beta|=k}[D^\beta u]_{\alpha;\Omega}.
\end{aligned}
\end{equation}
\begin{equation}
\tag{4.6}
\begin{aligned}
|u|_{k;\Omega}&=\sum_{j=0}^k[u]_{j,0;\Omega},
& |u|_{k,\alpha;\Omega}&=|u|_{k;\Omega}+[u]_{k,\alpha;\Omega}.
\end{aligned}
\end{equation}
If $d=\diam\Omega$, the dimensionless norms are
\[
|u|'_{k;\Omega}=\sum_{j=0}^k d^j[u]_{j,0;\Omega},
\qquad
|u|'_{k,\alpha;\Omega}=|u|'_{k;\Omega}
+d^{k+\alpha}[u]_{k,\alpha;\Omega}.
\tag{4.6'}
\]
These are Banach norms. If
$u\in C^\alpha$, $v\in C^\beta$, and $\gamma=\min(\alpha,\beta)$, then
\begin{equation}
\tag{4.7}
|uv|'_{0,\gamma;\Omega}
\leq |u|'_{0,\alpha;\Omega}|v|'_{0,\beta;\Omega}.
\end{equation}
\end{definition}

\section*{4.2. The Dirichlet Problem for Poisson's Equation}

\begin{lemma}\label{gt:lem-newtonian-c1}\dcref{ch4:4.1,ch4:4.8}\group{ch04-poisson-newtonian-potential}
\source{gilbarg-trudinger:page-0066}
If $f$ is bounded and integrable on $\Omega$, its Newtonian potential belongs
to $C^1(\mathbb R^n)$ and
\begin{equation}
\tag{4.8}
D_iw(x)=\int_\Omega D_i\Gamma(x-y)f(y)\,dy.
\end{equation}
\end{lemma}
\begin{proof}
\sketch Cut off the kernel on $|x-y|<\varepsilon$. Differentiation is valid
for the regularized potential, while the bounds
$|D\Gamma(z)|\leq C|z|^{1-n}$ and
$|\Gamma(z)|\leq C(1+|z|^{2-n}+|\log|z||)$ make both the potentials and
their first derivatives converge uniformly on compact sets.
\end{proof}

\begin{lemma}\label{gt:lem-newtonian-c2}\dcref{ch4:4.2,ch4:4.9}\group{ch04-poisson-newtonian-potential}
\source{gilbarg-trudinger:page-0067, gilbarg-trudinger:page-0068}
If $f$ is bounded and locally H\"older continuous on $\Omega$, then its
Newtonian potential is in $C^2(\Omega)$ and $\Delta w=f$. For any
divergence-theorem domain $\Omega_0\supset\Omega$, after extending $f$ by
zero,
\begin{equation}
\tag{4.9}
D_{ij}w(x)=\int_{\Omega_0}D_{ij}\Gamma(x-y)(f(y)-f(x))\,dy
-f(x)\int_{\partial\Omega_0}D_i\Gamma(x-y)\nu_j(y)\,ds_y.
\end{equation}
\end{lemma}
\begin{proof}
\sketch Differentiate the cut-off expression from Lemma~\ref{gt:lem-newtonian-c1},
subtract $f(x)$ to cancel the singularity, and integrate the constant part by
parts. H\"older continuity makes the remaining error $O(\varepsilon^\alpha)$.
Taking the trace in (4.9), with $\Omega_0$ a centered ball, gives
$\Delta w=f$.
\end{proof}

\begin{theorem}\label{gt:thm-dirichlet-poisson}\dcref{ch4:4.3,ch4:4.10}\group{ch04-poisson-newtonian-potential}
\source{gilbarg-trudinger:page-0068}
Let $\Omega$ be bounded and every boundary point regular for the Laplacian.
If $f$ is bounded and locally H\"older in $\Omega$, then
\[
\Delta u=f\quad\hbox{in }\Omega,
\qquad u=\varphi\quad\hbox{on }\partial\Omega
\]
has a unique classical solution for each continuous $\varphi$. On a ball
$B$, this solution is
\begin{equation}
\tag{4.10}
u(x)=\int_{\partial B}K(x,y)\varphi(y)\,ds_y
+\int_BG(x,y)f(y)\,dy.
\end{equation}
\end{theorem}
\begin{proof}
\sketch If $w$ is the Newtonian potential of $f$, solve the harmonic
Dirichlet problem with boundary value $\varphi-w$. Then $u=v+w$; uniqueness
is the maximum principle.
\end{proof}

\section*{4.3. Second-Derivative H\"older Estimates}

\begin{lemma}\label{gt:lem-newtonian-holder-ball}\dcref{ch4:4.4,ch4:4.11,ch4:4.12,ch4:4.13}\group{ch04-poisson-newtonian-potential}
\source{gilbarg-trudinger:page-0069, gilbarg-trudinger:page-0070, gilbarg-trudinger:page-0071}
\source{gilbarg-trudinger:page-0070}
\source{gilbarg-trudinger:page-0071}
Let $B_1=B_R(x_0)$ and $B_2=B_{2R}(x_0)$. If
$f\in C^\alpha(\overline B_2)$ and $w$ is its Newtonian potential over
$B_2$, then $w\in C^{2,\alpha}(\overline B_1)$ and
\begin{equation}
\tag{4.11}
|D^2w|'_{0,\alpha;B_1}\leq C|f|'_{0,\alpha;B_2},
\qquad C=C(n,\alpha).
\end{equation}
\end{lemma}
\begin{proof}
\sketch Formula (4.9) and the kernel estimates first give
\begin{equation}
\tag{4.12}
|D^2w(x)|\leq C\bigl(|f(x)|+R^\alpha[f]_{\alpha;B_2}\bigr).
\end{equation}
For $x,\bar x\in B_1$, split the difference of the two formulas into the
ball of radius $|x-\bar x|$ about their midpoint and its complement. H\"older
control handles the near integrals; the mean-value estimate for
$D^2\Gamma$ and integration by parts handle the far and boundary terms.
Thus
\begin{equation}
\tag{4.13}
|D^2w(x)-D^2w(\bar x)|
\leq C\bigl(R^{-\alpha}|f|_0+[f]_\alpha\bigr)|x-\bar x|^\alpha.
\end{equation}
\end{proof}

\begin{theorem}\label{gt:thm-poisson-compact-support}\dcref{ch4:4.5,ch4:4.14,ch4:4.15}\group{ch04-poisson-newtonian-potential}
\source{gilbarg-trudinger:page-0072}
If $u\in C_0^2(\mathbb R^n)$, $f\in C_0^\alpha(\mathbb R^n)$, and
$\Delta u=f$, then $u\in C_0^{2,\alpha}$ and, on any ball $B_R$ containing
the support,
\begin{equation}
\tag{4.14}
|D^2u|'_{0,\alpha;B_R}\leq C|f|'_{0,\alpha;B_R},
\qquad |u|_{1;B_R}\leq CR^2|f|_{0;B_R}.
\end{equation}
\end{theorem}
\begin{proof}
\sketch Green's formula reduces to
\begin{equation}
\tag{4.15}
u(x)=\int_{\mathbb R^n}\Gamma(x-y)f(y)\,dy.
\end{equation}
Apply Lemmas~\ref{gt:lem-newtonian-c1} and
\ref{gt:lem-newtonian-holder-ball}; the zeroth-order bound follows from the
first derivative bound and compact support.
\end{proof}

\begin{theorem}[Interior Schauder estimate for Poisson's equation]\label{gt:thm-poisson-interior-holder}\dcref{ch4:4.6,ch4:4.16}\group{ch04-poisson-newtonian-potential}
\source{gilbarg-trudinger:page-0072, gilbarg-trudinger:page-0073}
If $\Delta u=f$ in $\Omega$, with $u\in C^2(\Omega)$ and
$f\in C^\alpha(\Omega)$, then $u\in C^{2,\alpha}(\Omega)$. For concentric
$B_R\subset B_{2R}\Subset\Omega$,
\begin{equation}
\tag{4.16}
|u|'_{2,\alpha;B_R}
\leq C\bigl(|u|_{0;B_{2R}}+R^2|f|'_{0,\alpha;B_{2R}}\bigr),
\qquad C=C(n,\alpha).
\end{equation}
\end{theorem}
\begin{proof}
\sketch On $B_{2R}$ write $u=v+w$, where $v$ is harmonic and $w$ is the
Newtonian potential of $f$. Combine the harmonic derivative estimates with
Lemma~\ref{gt:lem-newtonian-holder-ball}. In dimension two, embed the
functions as coordinate-independent functions in dimension three for the
zeroth-order potential bound.
\end{proof}

\begin{corollary}\label{gt:cor-poisson-compactness}\dcref{ch4:4.7}\group{ch04-poisson-newtonian-potential}
\source{gilbarg-trudinger:page-0073}
Any bounded sequence of solutions of $\Delta u=f$ in $\Omega$, for fixed
$f\in C^\alpha(\Omega)$, has a subsequence converging uniformly on compact
subdomains to a solution.
\end{corollary}
\begin{proof}
\sketch Apply (4.16) on an exhaustion and use Arzela--Ascoli and a diagonal
subsequence.
\end{proof}

For $d_x=\dist(x,\partial\Omega)$ and
$d_{x,y}=\min(d_x,d_y)$, define the weighted interior norms
\begin{equation}
\tag{4.17}
\begin{aligned}
[u]^*_{k;\Omega}&=\sup_{x,|\beta|=k}d_x^k|D^\beta u(x)|,\\
[u]^*_{k,\alpha;\Omega}
&=\sup_{x\neq y,|\beta|=k}
d_{x,y}^{k+\alpha}
\frac{|D^\beta u(x)-D^\beta u(y)|}{|x-y|^\alpha},\\
|u|^*_{k,\alpha;\Omega}&=\sum_{j=0}^k[u]^*_{j;\Omega}
+[u]^*_{k,\alpha;\Omega}.
\end{aligned}
\end{equation}
\begin{equation}
\tag{4.18}
|f|^{(k)}_{0,\alpha;\Omega}
=\sup_x d_x^k|f(x)|+\sup_{x\neq y}d_{x,y}^{k+\alpha}
\frac{|f(x)-f(y)|}{|x-y|^\alpha}.
\end{equation}
If $\Omega$ is bounded with $d=\diam\Omega$, then
\begin{equation}
\tag{4.17'}
|u|^*_{k;\Omega}\leq\max(1,d^{k+\alpha})|u|_{k;\Omega}.
\end{equation}
If $\Omega'\Subset\Omega$ and $\sigma=\dist(\Omega',\partial\Omega)$, then
\begin{equation}
\tag{4.17''}
\min(1,\sigma^{k+\alpha})|u|_{k,\alpha;\Omega}
\leq|u|^*_{k,\alpha;\Omega}.
\end{equation}

\begin{theorem}[Weighted interior estimate]\label{gt:thm-poisson-interior-weighted}\dcref{ch4:4.8,ch4:4.19,ch4:4.20}\group{ch04-poisson-newtonian-potential}
\source{gilbarg-trudinger:page-0074}
If $\Delta u=f$ in an open $\Omega$, then
\begin{equation}
\tag{4.19}
|u|^*_{2,\alpha;\Omega}
\leq C\bigl(|u|_{0;\Omega}+|f|^{(2)}_{0,\alpha;\Omega}\bigr),
\qquad C=C(n,\alpha).
\end{equation}
\end{theorem}
\begin{proof}
\sketch Apply (4.16) on
$B_{d_x/3}(x)\subset B_{2d_x/3}(x)$ to obtain the weighted derivative bound
\begin{equation}
\tag{4.20}
|u|^*_{2;\Omega}\leq C(|u|_0+|f|^{(2)}_{0,\alpha}).
\end{equation}
For the H\"older term, use the same ball when $|x-y|<d_x/3$ and (4.20)
when the points are farther apart.
\end{proof}

\begin{theorem}[Unbounded forcing on a ball]\label{gt:thm-poisson-ball-unbounded}\dcref{ch4:4.9,ch4:4.21,ch4:4.22}\group{ch04-poisson-newtonian-potential}
\source{gilbarg-trudinger:page-0074, gilbarg-trudinger:page-0075}
Let $B$ be a ball, $f\in C^\alpha(B)$, and suppose for some
$0<\beta<1$ that
$\sup_Bd_x^{2-\beta}|f(x)|\leq N$. There is a unique
$u\in C^2(B)\cap C^0(\overline B)$ satisfying $\Delta u=f$ and
$u|_{\partial B}=0$, and
\begin{equation}
\tag{4.21}
\sup_Bd_x^{-\beta}|u(x)|\leq C(\beta)N.
\end{equation}
\end{theorem}
\begin{proof}
\sketch For $B_R(x_0)$, the barrier
$(R^2-|x-x_0|^2)^\beta$ gives
\begin{equation}
\tag{4.22}
|u(x)|\leq C N d_x^\beta.
\end{equation}
Truncate $f$, solve the bounded-data problems, and use
Corollary~\ref{gt:cor-poisson-compactness} on an exhaustion. The estimate
passes to the limit and gives the boundary values; uniqueness follows from
the maximum principle.
\end{proof}

\section*{4.4. Boundary Estimates}

Let $B_i^+=B_i\cap\mathbb R^n_+$ for concentric balls
$B_R\subset B_{2R}$, and let $T=\{x_n=0\}$.

\begin{lemma}\label{gt:lem-newtonian-holder-halfball}\dcref{ch4:4.10,ch4:4.23}\group{ch04-poisson-newtonian-potential}
\source{gilbarg-trudinger:page-0076}
If $f\in C^\alpha(\overline B_2^+)$ and $w$ is its Newtonian potential on
$B_2^+$, then $w\in C^{2,\alpha}(\overline B_1^+)$ and
\begin{equation}
\tag{4.23}
|D^2w|'_{0,\alpha;B_1^+}\leq C|f|'_{0,\alpha;B_2^+},
\qquad C=C(n,\alpha).
\end{equation}
\end{lemma}
\begin{proof}
\sketch Apply the proof of Lemma~\ref{gt:lem-newtonian-holder-ball} to
$B_2^+$. Boundary terms on $T$ vanish for tangential second derivatives;
recover $D_{nn}w$ from $\Delta w=f$.
\end{proof}

\begin{theorem}[Flat-boundary Schauder estimate]\label{gt:thm-poisson-halfball-holder}\dcref{ch4:4.11,ch4:4.24,ch4:4.25,ch4:4.26,ch4:4.27,ch4:4.28}\group{ch04-poisson-newtonian-potential}
\source{gilbarg-trudinger:page-0076, gilbarg-trudinger:page-0077}
\source{gilbarg-trudinger:page-0077}
\source{gilbarg-trudinger:page-0078}
If $\Delta u=f$ in $B_2^+$,
$u\in C^2(B_2^+)\cap C^0(\overline B_2^+)$,
$f\in C^\alpha(\overline B_2^+)$, and $u=0$ on $T$, then
$u\in C^{2,\alpha}(\overline B_1^+)$ and
\begin{equation}
\tag{4.24}
|u|'_{2,\alpha;B_1^+}
\leq C\bigl(|u|_{0;B_2^+}+R^2|f|'_{0,\alpha;B_2^+}\bigr).
\end{equation}
\end{theorem}
\begin{proof}
\sketch With $x^*=(x',-x_n)$, subtract the reflected Newtonian potential:
\begin{equation}
\tag{4.25}
w(x)=\int_{B_2^+}\bigl(\Gamma(x-y)-\Gamma(x^*-y)\bigr)f(y)\,dy.
\end{equation}
Even reflection of $f$, Lemmas~\ref{gt:lem-newtonian-holder-ball} and
\ref{gt:lem-newtonian-holder-halfball} give
\begin{equation}
\tag{4.26}
|D^2w|_{0,\alpha;B_1^+}\leq C|f|_{0,\alpha;B_2^+}.
\end{equation}
The difference $u-w$ is harmonic, vanishes on $T$, and extends oddly across
$T$, so the harmonic interior estimates prove (4.24). For compactly
supported $u$, the same representation yields
\begin{equation}
\tag{4.27}
|D^2u|_{0,\alpha;B_2^+}\leq C|f|_{0,\alpha;B_2^+},
\end{equation}
with
\begin{equation}
\tag{4.28}
u(x)=w(x)=\int_{B_2^+}
[\Gamma(x-y)-\Gamma(x^*-y)]f(y)\,dy.
\end{equation}
\end{proof}

For an open $\Omega\subset\mathbb R^n_+$ with flat boundary portion $T$,
define the partial interior norms as above using
$d_x=\dist(x,\partial\Omega\setminus T)$; in particular
\begin{equation}
\tag{4.29}
|u|^*_{k;\Omega\cup T}=\sum_{j=0}^k
\sup_{x,|\beta|=j}d_x^j|D^\beta u(x)|.
\end{equation}

\begin{theorem}[Weighted flat-boundary estimate]\label{gt:thm-poisson-partial-boundary}\dcref{ch4:4.12,ch4:4.29,ch4:4.30}\group{ch04-poisson-newtonian-potential}
\source{gilbarg-trudinger:page-0078}
If $\Delta u=f$ in $\Omega\subset\mathbb R^n_+$, $u=0$ on $T$, and
$f\in C^\alpha(\Omega\cup T)$, then
\begin{equation}
\tag{4.30}
|u|^*_{2,\alpha;\Omega\cup T}
\leq C\bigl(|u|_{0;\Omega}+|f|^{(2)}_{0,\alpha;\Omega\cup T}\bigr),
\qquad C=C(n,\alpha).
\end{equation}
\end{theorem}
\begin{proof}
\sketch Apply Theorem~\ref{gt:thm-poisson-halfball-holder} on half-balls of
radius comparable to $d_x$, and separate nearby from distant pairs as in
Theorem~\ref{gt:thm-poisson-interior-weighted}.
\end{proof}

\begin{theorem}[Regularity on a ball]\label{gt:thm-poisson-ball-boundary}\dcref{ch4:4.13,ch4:4.31,ch4:4.32}\group{ch04-poisson-newtonian-potential}
\source{gilbarg-trudinger:page-0079}
Let $B$ be a ball. If
$u\in C^2(B)\cap C^0(\overline B)$, $f\in C^\alpha(\overline B)$,
$\Delta u=f$, and $u=0$ on $\partial B$, then
$u\in C^{2,\alpha}(\overline B)$.
\end{theorem}
\begin{proof}
\sketch Translate a boundary point to the origin. Inversion and the Kelvin
transform
\begin{equation}
\tag{4.31}
v(x)=|x|^{2-n}u(x/|x|^2)
\end{equation}
map the ball locally to a half-space and satisfy
\begin{equation}
\tag{4.32}
\Delta v(x^*)=|x^*|^{-n-2}(\Delta u)(x).
\end{equation}
Apply Theorem~\ref{gt:thm-poisson-halfball-holder} at each boundary point
and the interior estimate away from the boundary.
\end{proof}

\begin{corollary}\label{gt:cor-poisson-ball-dirichlet}\dcref{ch4:4.14}\group{ch04-poisson-newtonian-potential}
\source{gilbarg-trudinger:page-0079}
If $\varphi\in C^{2,\alpha}(\overline B)$ and
$f\in C^\alpha(\overline B)$, then $\Delta u=f$ in $B$,
$u=\varphi$ on $\partial B$, has a unique
$u\in C^{2,\alpha}(\overline B)$.
\end{corollary}
\begin{proof}
\sketch Apply Theorems~\ref{gt:thm-dirichlet-poisson} and
\ref{gt:thm-poisson-ball-boundary} to $u-\varphi$.
\end{proof}

\section*{4.5. First-Derivative H\"older Estimates}

For
\begin{equation}
\tag{4.33}
\Delta u=\operatorname{div}f=D_if^i,
\end{equation}
use the generalized Newtonian potential
\begin{equation}
\tag{4.34}
w(x)=D_j\int_\Omega\Gamma(x-y)f^j(y)\,dy.
\end{equation}
If $f$ is H\"older continuous, subtracting $f(x)$ under the differentiated
singular integral gives
\begin{equation}
\tag{4.35}
D_iw(x)=\int_\Omega D_{ij}\Gamma(x-y)(f^j(y)-f^j(x))\,dy
-f^j(x)\int_{\partial\Omega}D_i\Gamma(x-y)\nu_j\,ds.
\end{equation}

\begin{theorem}[Interior $C^{1,\alpha}$ estimate]\label{gt:thm-poisson-first-derivative-interior}\dcref{ch4:4.15,ch4:4.33,ch4:4.34,ch4:4.35,ch4:4.36,ch4:4.37}\group{ch04-poisson-newtonian-potential}
\source{gilbarg-trudinger:page-0080}
\source{gilbarg-trudinger:page-0081}
If $u$ solves (4.33) in $\Omega$ with $f\in C^\alpha(\Omega)$, then for
concentric $B_R\subset B_{2R}\Subset\Omega$,
\begin{equation}
\tag{4.37}
|u|'_{1,\alpha;B_R}
\leq C\bigl(|u|_{0;B_{2R}}+R|f|'_{0,\alpha;B_{2R}}\bigr),
\qquad C=C(n,\alpha).
\end{equation}
\end{theorem}
\begin{proof}
\sketch The same near/far decomposition as in Lemma~\ref{gt:lem-newtonian-holder-ball}
applied to (4.35) gives
\begin{equation}
\tag{4.36}
|Dw|'_{0,\alpha;B_R}\leq C|f|'_{0,\alpha;B_{2R}}.
\end{equation}
Subtract $w$ from $u$ and apply the harmonic interior estimates.
\end{proof}

For the half-ball potential (4.34), the corresponding boundary estimate is
\begin{equation}
\tag{4.38}
|Dw|'_{0,\alpha;B_1^+}\leq C|f|'_{0,\alpha;B_2^+},
\qquad C=C(n,\alpha).
\end{equation}

\begin{theorem}[Flat-boundary $C^{1,\alpha}$ estimate]\label{gt:thm-poisson-first-derivative-halfball}\dcref{ch4:4.16,ch4:4.39,ch4:4.40,ch4:4.41,ch4:4.42,ch4:4.43,ch4:4.44,ch4:4.45,ch4:4.46}\group{ch04-poisson-newtonian-potential}
\source{gilbarg-trudinger:page-0081, gilbarg-trudinger:page-0082}
\source{gilbarg-trudinger:page-0082}
Let $u\in C^0(\overline B_2^+)$ solve (4.33) with
$f\in C^\alpha(\overline B_2^+)$ and $u=0$ on $T$. Then
\begin{equation}
\tag{4.43}
|u|'_{1,\alpha;B_1^+}
\leq C\bigl(|u|_{0;B_2^+}+R|f|'_{0,\alpha;B_2^+}\bigr).
\end{equation}
More generally, if
\begin{equation}
\tag{4.44}
\Delta u=g+\operatorname{div}f,
\end{equation}
then the corresponding interior and flat-boundary bounds are
\begin{equation}
\tag{4.45}
|u|'_{1,\alpha;B_1}
\leq C\bigl(|u|_{0;B_2}+R^2|g|_{0;B_2}
+R|f|'_{0,\alpha;B_2}\bigr).
\end{equation}
\begin{equation}
\tag{4.46}
|u|'_{1,\alpha;B_1^+}
\leq C\bigl(|u|_{0;B_2^+}+R^2|g|_{0;B_2^+}
+R|f|'_{0,\alpha;B_2^+}\bigr).
\end{equation}
\end{theorem}
\begin{proof}
\sketch Differentiate the half-space Green potential
\begin{equation}
\tag{4.39}
v(x)=-\int_{B_2^+}D_y
\bigl(\Gamma(x-y)-\Gamma(x-y^*)\bigr)\mathbin{\cdot}f(y)\,dy.
\end{equation}
Even reflection expresses its tangential and normal pieces as first
derivatives of full-ball Newtonian potentials. For $i<n$ and $i=n$,
respectively,
\begin{equation}
\tag{4.40}
v_i(x)=D_i\left[2\int_{B_2^+}\Gamma(x-y)f^i(y)\,dy
-\int_{B_2^+\cup B_2^-}\Gamma(x-y)f^i(y)\,dy\right],
\end{equation}
\begin{equation}
\tag{4.41}
v_n(x)=D_n\int_{B_2^+\cup B_2^-}\Gamma(x-y)f^n(y)\,dy.
\end{equation}
Equations (4.36) and (4.38) give
\begin{equation}
\tag{4.42}
|Dv|'_{0,\alpha;B_1^+}\leq C|f|'_{0,\alpha;B_2^+}.
\end{equation}
The remainder is
harmonic and vanishes on $T$. For (4.45)--(4.46), add the first-derivative
estimate for the Newtonian potential of bounded $g$.
\end{proof}
